% Ad Familiares 11.3 --- headnote
% ad-familiares-11-03, 4 August 44 BC, Naples

\textit{Marcus Brutus and Cassius, both praetors, to the
consul Mark Antony, from Naples on 4 August 44 BC --- Perseus
dateline \textit{Scr. Neapoli prid. Non. Sext. a. 710 (44)},
which the closing line of the letter itself confirms.
The salutation uses the formal opening \textit{si vales, bene
est}, a courtesy entirely at odds with what follows. The
addressee is Antony at Rome; the writers are by now on the
Campanian coast, having been forced from the city by the
veterans the previous letter (11.2) had asked him to control.}

\textit{This is the famous indignant reply --- the public
letter in which the surviving conspirators come out from
under the language of conciliation and answer Antony as
equals to an aggressor. He had issued an edict and a letter
of his own, ``abusive, threatening,'' which evidently threw
Caesar's death in their faces; they answer point by point.
The list in section 2 (levies, requisitions, tampered armies,
agents sent overseas) is a catalogue of Antony's actions over
the summer, here treated as already accomplished facts that
do not need denying. The third section refuses to be
intimidated --- ``Antony has no business commanding the men
by whose work he is a free man'' --- and the closing of
section 4 is the line most often quoted from the
correspondence of this year: do not consider how long Caesar
lived, but how short a time he reigned. Within weeks Cicero
would deliver the First Philippic, and the open contest with
Antony begin in earnest.}
